\documentclass[a4paper]{article}
\usepackage[margin=0.5in]{geometry}
\usepackage[fontsize=12pt]{scrextend}
\usepackage[english]{babel}
\usepackage{graphicx}
\usepackage{caption}
\usepackage{amsmath, amssymb}
\usepackage{algorithm}
\usepackage{amsmath}
\usepackage[noend]{algpseudocode}

\setlength{\parindent}{0em}
\setlength{\parskip}{0em}
\renewcommand{\baselinestretch}{1.5}

\title{4.1: The Maximum-Subarray Problem}
\author{}
\date{}

\begin{document}
\maketitle

\textit{\textbf{4.1-1}} \\
What does \texttt{FIND-MAXIMUM-SUBARRAY} return when all elements of $A$ are negative?\\
\textit{\textbf{Ans 4.1-1}} \\
It will return the greatest element of $A$.\\

\textit{\textbf{4.1-2}} Write pseudocode for the brute-force method of solving the maximum-subarray problem. Your procedure should run in $\Theta(n^2)$ time.

\textit{\textbf{Ans 4.1-2}} \\
\begin{algorithm}
	\caption{BRUTE-FORCE-FIND-MAXIMUM-SUBARRAY}
	\begin{algorithmic}[1]
		\Procedure{maxSubarray}{$arr, n$}
		\Statex
		\State $max-sum := -\infty$
		\State $low := 0$
		\State $high := 0$
		\For{$i := 1 \textbf{ to } n$}
			\State $sum := 0$
			\For{$j := i \textbf{ to } n$}
				\State $sum := sum + arr[j]$
				\If{$sum >  max-sum$}
					\State $max-sum := sum$
					\State $low := i$
					\State $high := j$
				\EndIf
			\EndFor
		\EndFor
		\State \textbf{return} $\left( low, high, max-sum\right)$
		\EndProcedure
	\end{algorithmic}
\end{algorithm}
\newpage
\textit{\textbf{4.1-3}} Implement both the brute-force and recursive algorithms for the maximum-subarray problem on your own computer. What problem size $n_0$ gives the crossover point at which the recursive algorithm beats the brute-force algorithm? Then, change the base case of the recursive algorithm to use the brute-force algorithm whenever the problem size is less than $n_0$. Does that change the crossover point? \\
\textit{\textbf{Ans 4.1-3}} On my computer:
\begin{enumerate}
 \item Brute Force Algorithm: 0.08s
 \item Divide and Conquer: 0.09s
\end{enumerate}

\textit{\textbf{4.1-4}} Suppose we change the definition of the maximum-subarray problem to allow the result to be an empty subarray, where the sum of the values of an empty subarray is $0$. How would you change any of the algorithms that do not allow empty subarrays to permit an empty subarray to be the result? \\
\textit{\textbf{Ans 4.1-4}} If the original algorithm returns a negative sum, returning an empty subarray instead. \\

\textit{\textbf{4.1-5}} Use the following ideas to develop a nonrecursive, linear-time algorithm for the maximum-subarray problem. Start at the left end of the array, and progress toward the right, keeping track of the maximum subarray seen so far. Knowing a maximum subarray $A[1..j]$, extend the answer to find a maximum subarray ending at index $j + 1$ by using the following observation: a maximum subarray $A[i..j + 1]$, is either a maximum subarray of $A[1..j]$ or a subarray $A[i..j + 1]$, for some $1 \le i \le j + 1$. Determine a maximum subarray of the form $A[i..j + 1]$ in constant time based on knowing a maximum subarray ending at index $j$.
\textit{\textbf{Ans 4.1-5}} \\
\begin{algorithm}
\caption{ITERATIVE-MAX-SUBARRAY}
\begin{algorithmic}[1]
\Procedure{iterativeFindMaxSubarray}{$arr, n$}
	\State $max-sum := -\infty$
	\State $sum := -\infty$
	\For{$i := 1 \textbf{ to } n$}
		\State $currHigh := i$
		\If{$sum > 0$}
			\State $sum := sum + arr[i]$
		\Else
			\State $currLow = i$
			\State $sum := arr[i]$
		\EndIf
		\If{$sum > max-sum$}
			\State $max-sum := sum$
			\State $low := currLow$
			\State $high := currHigh$
		\EndIf
	\EndFor
	\State \textbf{return} $\left( low, high, max-sum\right)$
\EndProcedure
\end{algorithmic}
\end{algorithm}
\end{document}